% Part: lambda-calculus
% Chapter: syntax
% Section: terms

\documentclass[../../../include/open-logic-section]{subfiles}

\begin{document}

\olfileid[fr]{lam}{syn}{trm}
\olsection{Termes}

Les termes du lambda-calcul se construisent
par induction à partir d'une réserve infinie
de variables $\Obj{v_0}$, $\Obj{v_1}$, \dots,
du symbole « $\lambd$ » et des parenthèses.
Nous emploierons $x$, $y$, $z$, \dots{} pour
désigner des variables, et $M$, $N$, $P$,
\dots{} pour désigner des termes.

\begin{defn}[Termes] \ollabel{defn:term}
L'ensemble des \emph{termes} du lambda-calcul
est défini inductivement par les règles suivantes :
\begin{enumerate}
  \item \ollabel{defn:term-var} Si $x$ est une
    variable, alors $x$ est un terme.
  \item \ollabel{defn:term-abs} Si $x$ est une
    variable et $M$ un terme, alors
    $(\lambd[x][M])$ est un terme.
  \item \ollabel{defn:term-app} Si $M$ et $N$
    sont des termes, alors $(MN)$ est un terme.
\end{enumerate}
\end{defn}

Un terme $(\lambd[x][M])$ obtenu selon
\olref{defn:term-abs} résulte d'une
\emph{abstraction}; le~$x$ qui suit le
symbole~$\lambd$ est appelé son
\emph{paramètre}. Un terme $(MN)$ obtenu
selon \olref{defn:term-app} résulte d'une
\emph{application}.

Les termes ainsi définis sont entièrement
parenthésés. Cette écriture peut devenir
pesante, comme le montre le terme
$(\lambd[x][((\lambd[x][x])(\lambd[x][(xx)]))])$.
Nous adopterons plus loin des conventions
permettant de supprimer certaines parenthèses.
La définition officielle permet néanmoins
de déterminer comment un terme est construit
selon la \olref{defn:term}. Ainsi, la dernière
étape de la formation du terme
$(\lambd[x][((\lambd[x][x])(\lambd[x][(xx)]))])$
est nécessairement une abstraction de
paramètre~$x$. Ce terme résulte de l'abstraction
du terme $((\lambd[x][x])(\lambd[x][(xx)]))$,
qui est l'application de deux termes. Chacun
de ces deux termes résulte lui-même d'une
abstraction, et ainsi de suite.

\begin{prob}
Décrire la formation du terme
$(\lambd[g][((\lambd[x][(g (x x))])
  (\lambd[x][(g (x x))]))])$.
\end{prob}
\begin{explain}
Dans cet exercice, les parenthèses délimitent chaque application
et chaque abstraction, suivant les règles de formation.
\end{explain}
\end{document}
